\section{实施资格是乘法闸门，不是交易清单}

本报告不提供下单价格、数量、建仓日期或分批计划。未来每一项具体候选仅在以下乘法闸门全部为1时才达到“可供独立组合审议”的状态：
\[
q_{effective,i}=q_{portfolio}\times q_{risk\ state}\times q_{research,i}
\times q_{valuation,i}\times q_{market,i}\times q_{concentration,i}
\times q_{account,i}.
\]
任一项为0，结果即为0。目前T3为空、$q_{portfolio}=0$，因此所有市场候选$q=0$。开放研究资格、历史分派、产品标签、单个公告或部分同日行情均不推定其他项通过。

\begin{exhibitbox}[不同结构的候选必须补齐的当前数据包]
\small
\begin{tabularx}{\exhibitboxwidth}{L{2.65cm}P P}
\toprule
\textbf{候选结构} & \textbf{单项研究与当前实施证据} & \textbf{失败路径} \\
\midrule
直接A股股票 & 公司/行业和治理档案、审计财务、收入至EPS桥、当前价格熊/基/牛估值、股利/自由现金流覆盖、催化与证伪、自由流通量、停复牌/涨跌停、价差、深度和容量。 & 未完整通过时保持在T1/T2；全部单项及组合闸门通过后才可进入T3/T4审议。 \\
ETF/基金 & 当前文件、指数/授权规则、持仓与重叠、当前价格/IOPV或参考值、跟踪和分配、重大证券借贷/衍生品、价差、深度、容量、最小单位与账户。 & 候选$q=0$；不得因标签或指数股息率自动替换。 \\
公募REIT & 项目、租户/承购方、运营KPI、TTM已结算现金、NAV日期与项目事件、债务/资本开支、税费、相关性、价差、深度和容量。 & 候选$q=0$；不得以高历史分派率强行填满额度。 \\
主权债/现金工具 & 发行人或保护机制、期限、到期/提前支取条款、现金流、税费、账户及集中度。 & 不作为已满足5\%总回报目标的假设。 \\
\bottomrule
\end{tabularx}
\end{exhibitbox}

\section{动态市场质量和容量约束}

交易价格相对适用参考值的偏离必须低于候选估值容忍度与20日历史保护线中的更严格者；价差必须低于全成本预算与20日相对价差保护线中的更严格者。拟议金额不得超过目标缺口、当日成交额、20日中位成交额和可执行深度所导出的最严格容量上限。每个输入必须同币种、为正且可追溯。REIT披露NAV不是实时IOPV，直接股票也不能以一个收盘价替代完整的停牌、涨跌停和容量评估。

缺失、过期、停牌、异常估值偏离、无可执行深度、账户税费未对账或任何集中度超限均令相应$q=0$。未通过的候选不获得市场配置权限；它可以在下一轮开放研究中与新的合规候选比较，但不能因为历史分配、回报目标或产品标签自动转入任何高收益、低信用、相关ETF、直接股票或“补仓”批次。

\section{从T3到T4的独立审批记录}

每一项拟进入组合的候选须保留一份带日期的记录：来源定位与时间戳、原始记录哈希、研究档案版本、估值和熊/基/牛假设、税后费后现金输入、流动性/容量计算、共同风险映射、拟议边际权重、组合压力变化、有效$q$、独立数据/估值/风险/运营复核身份以及失败时的现金回退位置。直接股票还须在T2/T3清单中列明代码，并与当前估值行和券商／一致预期材料逐一匹配；声明的T3股票数量必须等于该匹配清单数量。

审批顺序是：数据核验、单项研究复核、独立风险/集中度复核、投资组合治理与适当性复核，再到组合所有者确认是否进入审议。最后一步不授权订单、不指定经纪商、不设定实时价格，也不创建自动交易。

\section{现金流记录与投资者自主用途}

\begin{enumerate}
  \item 记录每笔\textbf{已结算}现金分配的日期、来源、总额、税、费用、净额、结算证据和返还性成分标记；
  \item 按预先冻结的台账规则单列返还性成分；无论处理方式如何，均不得称为收益或重复计入总回报；
  \item 收款后由投资者按自身账户、税务、流动性和授权条件决定保留、取出或不作处理；本报告不规定默认用途；
  \item 只有在投资者独立考虑新的配置时，才可在定期评审中比较有效$q=1$的T3候选；收款本身不触发自动配置；
  \item 卖出资产所得属于处置或再配置，不得伪装为分配、独立回报或目标达成。
\end{enumerate}

\section{今天与未来的边界}

今天没有市场交易动作。完成资料包只会形成一份可由独立风险与投资组合负责人审议的资格记录，并不自动授权交易。若之后发生投资者自主的实际交易决定，仍须在当时用真实账户、市场状态、流动性、停复牌、涨跌停和税费条件作独立判断。
